\section*{Chapter \ref{ch:b1:structure-spectroscopy} --- Spectroscopy for Structure}

\begin{solution}{exo:b1:structure-spectroscopy:1}
$\varepsilon = A/(\ell c) = 0.64/(1.00 \times \num{2.0e-5}) =
\qty{3.2e4}{L/(mol.cm)}$. Three times the concentration: $A = 1.92$, if
the law still holds at that \omterm{def:b1:structure-spectroscopy:absorbance}{absorbance}.
\end{solution}

\begin{solution}{exo:b1:structure-spectroscopy:2}
\ce{C6H6}: $D = (12 + 2 - 6)/2 = 4$, benzene (a ring and three $\pi$ bonds).
\ce{C4H8O2}: 1, ethyl ethanoate or butanoic acid. \ce{C3H7NO}: $(6 + 2 + 1
- 7)/2 = 1$, propanamide. \ce{C2H3Cl}: 1, chloroethene. \ce{C10H14O}:
$(20 + 2 - 14)/2 = 4$: carvone has one ring, two C=C and one C=O.
\end{solution}

\begin{solution}{exo:b1:structure-spectroscopy:3}
Propanone: one signal (6H). Ethanoic acid: two (3H, 1H). 2-Methylpropane:
two (9H, 1H). 1,2-Dichloroethane: one (4H).
\end{solution}

\begin{solution}{exo:b1:structure-spectroscopy:4}
$\delta = 1460/400 = \qty{3.65}{ppm}$; at \qty{90}{MHz} the signal lies
$3.65 \times 90 = \qty{329}{Hz}$ from the reference. \omterm{def:b1:structure-spectroscopy:coupling}{Coupling constants}
stay the same in hertz while shift differences grow with the field:
\omterm{def:b1:structure-spectroscopy:coupling}{multiplets} that overlap at low field separate at high field.
\end{solution}

\begin{solution}{exo:b1:structure-spectroscopy:5}
\ce{CH3}: triplet (3H); central \ce{CH2}: five neighbours, sextet
1\,:\,5\,:\,10\,:\,10\,:\,5\,:\,1 (2H); \ce{CH2Br}: triplet (2H), the most
deshielded.
\end{solution}

\begin{solution}{exo:b1:structure-spectroscopy:6}
Ethanol: O--H (3666) and no C=O; propanone: C=O (1738) and no O--H;
ethanoic acid: both, O--H (3582) and C=O (1788).
\end{solution}

\begin{solution}{exo:b1:structure-spectroscopy:7}
Ethyl ethanoate: quartet (2H) near 4.1, singlet (3H) near 2.0, triplet
(3H) near 1.3 ppm. Methyl propanoate: singlet (3H) of \ce{OCH3}, deshielded
by the oxygen (near 3.7 ppm), quartet (2H) of \ce{CH2C=O} (near 2.3), triplet
(3H) near 1.1. The quartet is the most deshielded signal in the first and
not in the second: a quartet near 4 ppm means \ce{O-CH2CH3}.
\end{solution}

\begin{solution}{exo:b1:structure-spectroscopy:8}
Lines \qty{0.080}{ppm} apart, that is $0.080 \times 90 = \qty{7.2}{Hz}$:
the \omterm{def:b1:structure-spectroscopy:coupling}{coupling constant}. The triplet has the same spacing: the \ce{CH2} and
\ce{CH3} are coupled to each other.
\end{solution}

\begin{solution}{exo:b1:structure-spectroscopy:9}
A C=O band and a single signal: propanone, \ce{CH3COCH3}. The aldehyde
isomer, propanal \ce{CH3CH2CHO}, shows the C--H stretch of the
\ce{CHO} group and its very deshielded proton.
\end{solution}

\begin{solution}{exo:b1:structure-spectroscopy:10}
The first set splits the line into $n_1 + 1$ lines; each of these is split
by the second set into $n_2 + 1$: $(n_1 + 1)(n_2 + 1)$ lines. If $J_1 =
J_2$, a line's position depends only on $k_1 + k_2$, which takes $n_1 + n_2 +
1$ values: the $n + 1$ rule with $n = n_1 + n_2$. Central \ce{CH2} of
1-bromopropane: $4 \times 3 = 12$ lines in general, six (a sextet) when the
two constants are equal.
\end{solution}

\begin{solution}{exo:b1:structure-spectroscopy:11}
$D = 0$; an alcohol (O--H band, no C=O). Doublet (6H): two equivalent
\ce{CH3} on one CH; \omterm{def:b1:structure-spectroscopy:coupling}{multiplet} (1H): that CH; doublet (2H): a \ce{CH2}
next to the CH and bearing the OH; broad singlet: OH.
2-Methylpropan-1-ol, \ce{(CH3)2CH-CH2-OH}.
\end{solution}

\begin{solution}{exo:b1:structure-spectroscopy:12}
Each species absorbs independently: $A = A_1 + A_2 = \ell(\varepsilon_1c_1
+ \varepsilon_2c_2)$. At two wavelengths, with the four coefficients known,
two linear equations give $c_1$ and $c_2$.
\end{solution}

\begin{solution}{pb:b1:structure-spectroscopy:1}
\textbf{1.} $0.5690 \times 12.0/44.0 = \qty{0.1552}{g}$.
\textbf{2.} $0.2328 \times 2.0/18.0 = \qty{0.02587}{g}$.
\textbf{3.} $0.2500 - 0.1552 - 0.0259 = \qty{0.0689}{g}$.
\textbf{4.} C \qty{62.1}{\percent}, H \qty{10.3}{\percent}, O
\qty{27.6}{\percent}.
\textbf{5.} C $0.1552/12.0 = 0.01293$, H $0.02587$, O $0.0689/16.0 =
0.00431$ mol: ratio $3.00 : 6.00 : 1$, \ce{C3H6O}.
\textbf{6.} $M = \rho RT/p = 3740 \times 8.314 \times 373.15/(1.000 \times 10^5) =
\qty{116}{g/mol}$ (density in \unit{g/m^3}).
\textbf{7.} $116/58 = 2$: \ce{C6H12O2}, $D = (12 + 2 - 12)/2 = 1$.
\textbf{8.} A C=O stretch.
\textbf{9.} Any O--H: no acid, no alcohol.
\textbf{10.} A C--O stretch, strong: the C--O of an ester.
\textbf{11.} C--H stretches of \ce{CH2} and \ce{CH3} groups.
\textbf{12.} One $\pi$ bond, used by the C=O; with a strong C--O and no
O--H, an ester. An acid needs an O--H; a diol has no C=O. A ketone
carrying an ether group would have its C=O near the \qty{1738}{cm^{-1}}
of propanone, while the observed \qty{1754}{cm^{-1}} lies close to the
\qty{1764}{cm^{-1}} of ethyl ethanoate; the NMR of Part III confirms the
ester.
\textbf{13.} Five; $2 + 2 + 2 + 3 + 3 = 12$, the twelve hydrogens.
\textbf{14.} A \ce{CH2} with three neighbours (a \ce{CH3}), bound to the
oxygen (deshielded): \ce{-O-CH2-CH3}.
\textbf{15.} The \ce{CH3} of that ethyl group, with two neighbours; the two
signals are coupled to each other, hence one $J$.
\textbf{16.} A \ce{CH2} with two neighbours, next to the C=O.
\textbf{17.} A \ce{CH2} with five neighbours: between a \ce{CH2} and a
\ce{CH3}.
\textbf{18.} A \ce{CH3} with two neighbours, far from the functional group.
\textbf{19.} \ce{CH3CH2O-} and \ce{-CH2CH2CH3}.
\textbf{20.} The 4.13 ppm protons are on the carbon bound to oxygen, the
2.28 ppm protons next to the C=O: both electron-withdrawing.
\textbf{21.} \ce{CH3CH2-O-CO-CH2CH2CH3} or \ce{CH3CH2-CO-O-CH2CH2CH3}. Only
the first puts the \ce{CH2} with three neighbours on the oxygen (the
quartet at 4.13 ppm).
\textbf{22.} Propyl propanoate is the second: its \ce{OCH2} would be a
triplet and the quartet would lie near 2.3 ppm.
\textbf{23.} Butyl ethanoate, \ce{CH3CO-O-C4H9}, would show a singlet (3H)
near 2.0 ppm and no quartet.
\textbf{24.} Ethyl butanoate, \ce{CH3CH2CH2-CO-O-CH2CH3}.
\textbf{25.} IR: C=O 1754, C--O 1182, C--H 2982, no O--H. NMR: \ce{OCH2}
4.13 (q), \ce{CH2CO} 2.28 (t), central \ce{CH2} 1.66 (sextet), ester
\ce{CH3} 1.25 (t), chain \ce{CH3} 0.95 (t), \omterm{def:b1:structure-spectroscopy:integration}{integrations} 2, 2, 2, 3, 3.
\textbf{26.} \textbf{\qty{116}{g/mol}}: ethyl butanoate, \ce{C6H12O2}.
\end{solution}
